% from ega4-16-fr.tex
\oldpage[IV-4]{16}
\begin{proposition}[16.3.8]
\label{IV.16.3.8}
L'image de l'homomorphisme $d_{X/S}^n:\mathcal{O}_X\to\mathscr{P}_{X/S}^n$
engendre le $\mathcal{O}_X$-Module $\mathscr{P}_{X/S}^n$.
\end{proposition}

On se ramène aussitôt au cas où $X=\operatorname{Spec}(B)$ et
$S=\operatorname{Spec}(A)$ sont affines et la proposition résulte de (16.3.7.1)
puisque $\pi_n$ est surjectif. On notera qu'en général $d_{X/S}^n$ n'est pas
surjectif (même déjà pour $n=1$).

\begin{proposition}[16.3.9]
\label{IV.16.3.9}
Supposons que $f:X\to S$ soit un morphisme localement de type fini. Alors les
$\mathscr{P}_f^n$ et les $\mathscr{G}r_n(\mathscr{P}_f)$ sont des
$\mathcal{O}_X$-Modules quasi-cohérents de type fini.
\end{proposition}

Cela résulte de (16.1.6) et du fait que $\Delta_f$ est localement de présentation
finie (I, 4.3.1).

\subsection{Propriétés fonctorielles des invariants différentiels}
\EGAReaderTarget{ega4.c04.subsec.16-4}{IV.16.4}

\begin{env}[16.4.1]
\label{IV.16.4.1}
Considérons un diagramme commutatif de morphismes de préschémas
\[
  \xymatrix{
    X \ar[d]_-{f} & X' \ar[l]_-u \ar[d]^-{f'} \\
    S & S' \ar[l]^-w
  }
  \tag{16.4.1.1}
  \label{IV.16.4.1.1}
\]
On en déduit un diagramme commutatif
\[
  \xymatrix{
    X \ar[d]_-{\Delta_f} & X' \ar[l]_-u \ar[d]^-{\Delta_{f'}} \\
    X\times_S X & X'\times_{S'}X' \ar[l]^-v
  }
\]
où $v$ est l'homomorphisme composé (I, 5.3.5 et 5.3.15)
\[
  X'\times_{S'}X'
  \xrightarrow{(p'_1,p'_2)_S}X'\times_S X'
  \xrightarrow{u\times_S u}X\times_S X.
  \tag{16.4.1.2}
  \label{IV.16.4.1.2}
\]
On déduit donc de $u$ et $v$, comme il a été expliqué dans (16.2.1), des
homomorphismes de faisceaux d'anneaux augmentés
\[
  \nu_n:\rho^*(\mathscr{P}_{X/S}^n)\longrightarrow
  \mathscr{P}_{X'/S'}^n
  \tag{16.4.1.3}
  \label{IV.16.4.1.3}
\]
(où l'on a posé $u=(\rho,\lambda)$); ces homomorphismes forment un système
projectif, et donnent donc à la limite un homomorphisme de faisceaux d'anneaux
augmentés
\[
  \nu_\infty:\rho^*(\mathscr{P}_{X/S}^\infty)\longrightarrow
  \mathscr{P}_{X'/S'}^\infty;
  \tag{16.4.1.4}
  \label{IV.16.4.1.4}
\]
d'autre part, par passage aux quotients, les homomorphismes $\nu_n$ donnent un
di-homomorphisme d'Algèbres graduées (relatif à $\lambda^\#$):
\[
  \operatorname{gr}(u):\rho^*(\mathscr{G}r_\bullet(\mathscr{P}_{X/S}))
  \longrightarrow\mathscr{G}r_\bullet(\mathscr{P}_{X'/S'}).
  \tag{16.4.1.5}
  \label{IV.16.4.1.5}
\]
\end{env}

\begin{env}[16.4.2]
\label{IV.16.4.2}
